%Orbiter provides support for matrix groups and their various permutation representationes. 
Tables~\ref{tab:matrix:group:1}-\ref{tab:matrix:group:2} shows the Orbiter commands to create a matrix group.
\orbitertableofcommands{linear_group_1}
\orbitertableofcommands{linear_group_2}
Matrix groups are treated as permutation groups, 
so that a stabilizer chain can be established. 
Therefore, when creating a matrix group, a permutation representation must be chosen.
By default, this will be 
the natural permutation representation. 
In the case of a projective matrix group, this is the action on the points of the underlying projective space.
In case of an affine group or a general linear group, 
is is the set of vectors of the underlying vector space.
At times, there are other permutation actions which are more suitable. 
Here, suitable often means that the degree of the permutation representation is small. 
This helps improve the performance of the permutation group algorithms. 
For background information about the classical groups of matrices 
over finite fields and their actions, 
see~\cite{Taylor1992}, for instance. 


\bigskip

Let $V \simeq \bbF_q^n$
be a finite dimensional vector space over $\bbF_q$. 
The set of subspaces of $V$ form the projective geometry $\PG(n-1,q).$
The points of the geometry are the one-dimensional subspaces of $\bbF_q^n.$
The hyperplanes are the codimension one subspaces. 
For each $i$ with $0 \le i \le n$, there are subspaces of that dimension. 
They form the elements of a given type. 
For $i=1$ and $i=n-1$ we have points and hyperplanes. 
For $i=2$ we get lines. 
For $i=3$ we get planes.
The easiest way to denote a subspace if by means of a generator matrix. 
This is a matrix whose rows form a basis for the subspace in question. 


\bigskip

Let $\pi$ be a projective space.
A collineation of a projective space $\pi$ 
is a bijective mapping from the points of $\pi$ to themselves  
which preserves collinearity. That is, a collineation $\varphi$ maps any 
three collinear points $P,Q,R$ 
to another collinear triple $\varphi(P), \varphi(Q), \varphi(R).$ 
The collineations form a group with respect to composition, the collineation group.
If $M$ is the matrix of an endomorphism, then $\Psi_M$ is the induced 
map on projective space. 
By considering the homomorphism $M \mapsto \Psi_M$, 
the group $\GL(n+1,q)$ of invertible endomorphisms becomes 
a subgroup of the group of collineations of $\PG(n,q).$ 
This is the projectivity group $\PGL(n+1,q).$ 
It is isomorphic to $\GL(n+1,q)/\bbF_q^\times.$ 
Another source of collineations is this: 
Let $\Phi \in \Aut(\bbF_q)$ be a field automorphism. Then 
$\Phi$ acts on projective space by sending 
$P(\bx)$ to $P(\bx\Phi).$ 
This map is another type of collineation, 
called automorphic collineation.
This way, $\Aut(\bbF_q)$ gives rise to a group of collineations.
If $q = p^h$ for some prime $p$ and some integer $h$ then 
$$
\Phi_0: \bbF_q \rightarrow \bbF_q, \; x \mapsto x^p
$$ 
is a generator for the cyclic group $C_h \simeq \Aut(\bbF_q).$ 
The collineation group 
of $\PG(n,q)$ ($n \ge 2$) is isomorphic 
to the semidirect product of the projectivity group and the 
automorphism group of the field. 
The collineation group is 
$\PGGL(n+1,q) = \PGL(n+1,q) \ltimes \Aut(\bbF_q).$
We use the following notation for elements of $\PGGL(n+1,q).$ 
Let $\Phi_0$ be a generator 
for $\Aut(\bbF_q)$ and let $M \in \GL(n+1,q).$ 
The map 
$$
(\Psi_M,\Phi_0^k): \PG(n,q) \rightarrow \PG(n,q), \; 
P(\bx) \mapsto P(\by), \quad \by = (\bx \cdot M)^{\Phi_0^k}
$$
is denoted as 
\begin{equation}\label{eqn:Mk}
M_k.
\end{equation}
The identity element is $I_0$, 
where $I$ is the identity matrix and $0$ is the residue class modulo $h.$
The rules for multiplication and inversion in the collineation group are given as
\begin{eqnarray}\label{eqn:MkNl}
M_k \cdot N_l &=& \Big(M \cdot N^{\Phi^{-k}}\Big)_{k+l},\\
\Big(M_k\Big)^{-1} &=& \Big(\Big(M^{-1}\Big)^{\Phi^{k}}\Big)_{-k}.
\end{eqnarray}
The affine group $\AGL(n,q)$ is the semidirect product of $\GL(n,q)$ with $\bbF_q^n.$ 
The affine semilinear group $\AGGL(n,q)$  is the semidirect product of $\AGL(n,q)$ with $\Aut(\bbF_q).$
The elements of $\AGGL(n,q)$ are triples 
$$
M_{\ba,k} := (M, \ba, k) \in \GL(n,q) \times \bbF_q^n \times \Aut(\bbF_q), 
$$
which act on $\bbF_q^n$:
$$
\Big(\bx, (M, \ba, k)\Big) \mapsto \big(\bx \cdot M + \ba\big)^{\Phi^k}.
$$
The multiplication in $\AGGL(n,q)$ is 
$$
M_{\ba,k}  \cdot N_{\bb,l} = (MN)_{\ba N^{\Phi^{-k}}+ \bb^{\Phi^{-k}},k+l}. 
$$
The inverse of an element is 
$$
\Big(M_{\ba,k} \Big)^{-1} = \Big( M^{-1}\Big)_{\ba^{\Phi^{k}} M^{-1},-k}.
$$

 


\bigskip

A correlation is a one-to-one mapping between the set of 
points and the set of hyperplanes which reverses incidence. 
So, if $\rho$ is a correlation and 
$P$ is a point and $\ell$ is a hyperplane then $P^\rho$ is a hyperplane and $\ell^\rho$ is 
a point and 
$$
\ell^\rho \in P^\rho \iff P \in \ell.
$$
A correlation of order two is called polarity.
The standard polarity is the map 
$$
\rho:  \cP \leftrightarrow \cL, \;P(\bx) \leftrightarrow [\bx].
$$





\bigskip

A group $G$ can act on $V$ in one of the types listed in Table~\ref{tab:actions}.

\orbitertableofextracommands{basic_actions}


\bigskip

To create a matrix group over a finite field, 
it is possible to specify the field by its order. 
It is also possible to use a specific field in the construction of the group. 
In this case, one can create the field beforehand, 
and reference the field object in the creation of the group.
The second way allows more control, as it allows to select  
the irreducible polynomial that is used to create the field. 
If only the field order is given, the default polynomial is used. 



\bigskip

The command 

\sourcecode{GL_2_3_elements_dense}

creates the group $\GL(2,3)$ acting on the $9$ vectors of $\bbF_3^2$. 
In this case, the field is created explicitly. The group has 48 elements, which are printed in latex format. 
The latex output is shown below:

\orbiteroutput{GRAPHICS/LATEX/GL_2_3_elements_dense}

\bigskip




The command 

\sourcecode{PGL_2_3_elements_dense}

creates the group $\PGL(2,3)$ acting on the $4$ points of $\PG(1,3)$. 
In this case, the field is created implicitly. The group has 24 elements, shown below:


\orbiteroutput{GRAPHICS/LATEX/PGL_2_3_elements_dense}



\bigskip

The command 

\sourcecode{PGL_2_17}

creates the group $\PGL(2,17)$ of order $4896$ acting on the $18$ points of $\PG(1,18)$. 
The group is not simple. 
However, the subgroup $\PSL(2,17)$ of order $2448$ is. 
To create it, we used the following command:

\sourcecode{PSL_2_17}

\bigskip

The command 

\sourcecode{PGL_4_2}

creates the group $\PGL(4,2)$ of order $20160$ acting on the $15$ elements of $\PG(3,2)$. 
In this case, the field is created implicitly.



\bigskip


The command 

\sourcecode{PGL_4_2_with_field}

creates the field $\bbF_2$ first, and the builds the group based in this field. 
The group is the same as the group created by the previous command.
However, it is good to be able to have more options when creating the finite field. 
For instance, when using extension fields, it is sometimes necessary to pick the appropriate polynomial 
that defines the field, as described in Section~\ref{sec:extension:fields}.


\bigskip

The command 

\sourcecode{PGL_4_2_report}

creates a report about the group $\PGL(4,2)$.
The report gives information about the underlying field and the projective geometry, 
as well as about the permutation action, 
including the stabilizer chain that is used to represent the group effectively.

\orbiteroutput{GRAPHICS/LATEX/report_group_PGL_4_2}



\bigskip


We use the following Orbiter command creates $\PGL(4,2)$ again.
The command invokes two activities. 
The first creates a latex report for the group in the file \verb'PGL_4_2_report.tex'.
The second activity exports the permutation representation in Orbiter makefile format.

\sourcecode{PGL_4_2_export}

The file \verb'PGL_4_2.makefile' is created:

\sourcecode{PGL_4_2_generated}

This command can be used to recreate the group as permutation group directly.
This group will be considered again in Section~\ref{sec:lineargroups} below.
The permutation representation itself is stored 
in the file \verb'PGL_4_2_gens.csv':
\orbiteroutput{GRAPHICS/LATEX/PGL_4_2_gens_csv}

\bigskip



The command

\sourcecode{L_5_3}

creates $\PSL(5,3)$ of order $237783237120$ 
in the action on the 121 points of $\PG(4,3).$

\bigskip

The command

\sourcecode{L_4_5}

creates $\PSL(4,5)$ of order $7254000000$ 
in the action on the 156 points of $\PG(3,5).$

\bigskip


The command

\sourcecode{PGL_4_5}

creates $\PGL(4,5)$ of order $29016000000$ 
in the action on the 156 points of $\PG(3,5).$
The group $\PSL(4,5)$ is a subgroup of $\PGL(4,5)$ of index 4.

\bigskip



The command

\sourcecode{PGGL_3_4}

creates $\PGGL(3,4)$ of order $120960$ 
in the action on the 21 points of $\PG(2,4).$
The report contains one $p$-Sylow subgroup 
for each of the primes $p$ dividing the order of the group.

\bigskip



The command

\sourcecode{PGGL_3_8}

creates $\PGGL(3,8)$ of order $49448448$ 
in the action on the 73 points of $\PG(2,8).$
The report contains one $p$-Sylow subgroup for each 
of the primes $p$ dividing the order of the group.

\bigskip






The command

\sourcecode{AGL_1_27}

creates $\AGL(1,27)$ of order $702$ 
in the action on the 27 points of $\AG(1,27).$

\bigskip




The command

\sourcecode{AGGL_2_27}

creates $\AGGL(2,27)$ of order $1117679472$ 
in the action on the 729 points of $\AG(2,27).$

\bigskip







The command 

\sourcecode{SP_4_2}

creates the symplectic group $\Sp(4,2)$ of order $720$ 
acting on the 16 points of $\AG(4,2).$


\bigskip







The command 

\sourcecode{PSP_4_4}

creates the symplectic group $\PSp(4,4)$ of order $979200$ 
acting on the 85 points of $\PG(3,4).$


\bigskip



The command 

\sourcecode{PGO_5_2}

creates the group $\PGO(5,2)$ acting on the 15 points of the $Q(4,2)$ quadric. 
The following latex report is produced:
\orbiteroutput{GRAPHICS/LATEX/report_group_PGO_5_2}

The group order is $720$.


\bigskip



The command 

\sourcecode{PGGO_5_4}

creates the group $\PGGO(5,4)$ of order $1958400$ 
acting on the 85 points of the $Q(4,4)$ quadric. 


\bigskip



The command 

\sourcecode{PGOp_6_2}

creates the group $\PGGO^+(6,2)$ of order $40320$ 
acting on the 35 points of the $Q^+(5,2)$ quadric. 


\bigskip


The command 

\sourcecode{PGOm_6_2}

creates the group $\PGGO^-(6,2)$ of order $51840$ 
acting on the 27 points of the $Q^-(5,2)$ quadric. 


\bigskip





The symplectic group $\PSp(6,2)$ can be created using the following command: 

\sourcecode{PSP_6_2}

The group order is $1451520$. 
The degree of the action is $63$.


\bigskip

The group $\PGO(7,2)$, isomorphic to $\PSp(6,2)$, 
can be created using the following command: 

\sourcecode{PGO_7_2}

The group order is $1451520$. The degree of the action is $63$.


\bigskip

The command 

\sourcecode{NullPolarity_6_2}

creates the null polarity group of $\PG(6,2)$ 
acting on the 63 points of $\PG(5,2)$. 
This is the group preserving the 
dot product
$$
\bx \cdot \by = \sum_{i=1}^n x_i y_i.
$$
The group order is $23040$. 
The degree of the action is $63$.


\bigskip


The command 

\sourcecode{M24_matrix_group}

creates the matrix group $M_{24}$ as a subgroup of $\PGL(12,2).$ 
The group is given by means of a generating set of 12 matrices, 
stored in the makefile variable 
$$
\verb'GOLAY_24_CODE_AUT_GENS'.
$$




\bigskip



